\documentclass{article}

\title{Ch. 31 CW Solutions}
\author{Principles of Macroeconomics}
\date{\today}

\begin{document}

\maketitle

\begin{enumerate}

\item No, the rate of money growth should equal the rate of real GDP growth. Recall the quantity equation is \[MV = PY.\] If velocity is constant and the Fed wants to keep the price level constant, then the only two variables that change in the quantity equation are $M$ and $Y$. In order for the equation to hold with equality, $M$ and $Y$ will have to change at the same rate.

\item If inflation increases sharply, the inflation tax increases on holders of money because their money loses its value more quickly. Wealth held in savings accounts is not subject to a change in inflation because savings accounts pay nominal interest rates that compensate holders for inflation. Nevertheless, inflation hurts holders of savings accounts because inflation increases the tax burden on savings.

\item

	\begin{enumerate}
	
	\item Before-tax: $r = 10 - 5 = 5\%$ \\
		After-tax: $r = 10(0.6) - 5 = 1\%$
		
	\item Before-tax: $r = 6 - 2 = 4\%$ \\
		After-tax: $r = 6(0.6) - 2 = 1.6\%$
		
	\item Before-tax: $r = 4 - 1 = 3\%$ \\
		After-tax: $r = 4(0.6) -1 = 1.4\%$
	
	\end{enumerate}
	
\item Inflation doesn't really affect money's role as the medium of exchange unless hyperinflation makes people unwilling to accept money in exchange for a good/service.

Inflation makes money a less consistent unit of account because it changes the value of money. That's a minimal problem as long as inflation is low and stable. 

Inflation makes money a worse store of value because it depreciates money. 

\item

	\begin{enumerate}
	
	\item It helps the government because it increases the inflation tax. The government is able to fund itself with lower taxes.
	
	\item It helps a homeowner with a fixed-rate mortgage because unexpectedly high inflation will make the real value of their mortgage payments lower than expected. 
	
	\item It hurts a union worker with a long-term contract because unexpectedly high inflation will make the real value of their wages lower than expected.
	
	\item It hurts a college that has invested some of its endowment in government bonds. Unexpectedly high inflation will make the real value of the bonds' interest payments lower than expected.
	
	\end{enumerate}
	
\item

	\begin{enumerate}
	
	\item Uncertain -- actual inflation that matches expected inflation doesn't hurt or help borrowers and lenders because they will have agreed to a nominal interest rate that accounts for expected inflation. If actual inflation is less than expected inflation, it hurts borrowers and helps lenders because the real value of interest payments is higher than expected. If actual inflation is higher than expected, it helps borrowers and hurts lenders because the real value of interest payments is lower than expected. 
	
	\item False -- changes in relative prices help some people and hurt others.
	
	\item True -- generally, wages inflate at the same rate as prices.
	
	\end{enumerate}

\end{enumerate}

\end{document}